\documentclass[10pt,a4paper]{amsart}
\usepackage[margin=19mm]{geometry}
\usepackage{amsmath,amssymb,mathtools}
\usepackage[scr=rsfs,cal=euler]{mathalfa}
\usepackage{xfrac}
\usepackage[unicode,hidelinks,bookmarks=false]{hyperref}
\newcommand{\ie}{i.e.\ }
\newcommand{\cf}{cf.\ }
\newcommand{\resp}{respectively\ }
\newcommand{\Subsection}{\subsection}
\providecommand{\nobreakdash}{\nobreak\hbox{-}}
\setlength{\emergencystretch}{2em}
\multlinegap0pt
\def\T{{ \mathrm{\scriptscriptstyle T} }}
\theoremstyle{plain}
\newtheorem{theorem}{Theorem}
\newtheorem{proposition}[theorem]{Proposition}
\newtheorem{lemma}[theorem]{Lemma}
\newtheorem{corollary}[theorem]{Corollary}
\theoremstyle{definition}
\newtheorem{definition}[theorem]{Definition}
\newtheorem{assumption}[theorem]{Assumption}
\theoremstyle{remark}
\newtheorem{remark}[theorem]{Remark}
\title[Minimax Robust Designs for M-Estimated Models]{Minimax Robust Designs for M-Estimated Models — Lemma 3}
\author{Rui Hu, Douglas P. Wiens}
\date{Source: arXiv:2604.21998v3 (CC BY 4.0)}
\begin{document}
\maketitle

\noindent\emph{Source and licence.} The delivered work is the article shown in the title above, version arXiv:2604.21998v3 (CC BY 4.0), distributed under the Creative Commons Attribution 4.0 International licence, as recorded in the arXiv licence metadata for this exact version and shown on the abstract page saved with this item. The full licence notice, which carries the licence URL and the address of that abstract page, is the bundled file \texttt{licenses/arxiv-2604.21998-CC-BY-4.0.txt}.\par\smallskip

\noindent\textit{Editorial scope.} Counted unit: the article's Lemma 3 (source0.tex lines 665--690, source label \texttt{lemma: limit}); statement and proof are the article's own text line for line, with no line omitted, substituted or re-flowed. The counted proof is the article's own run-in proof block, delivered exactly from the heading 'Proof of Lemma 3' through the line before the next proof heading of the paper. Editorial formatting only: an AMS \texttt{amsart} preamble that restates the article's own macro definition of \texttt{\textbackslash T} (which the article defines after its \texttt{\textbackslash begin\{document\}}) and the theorem environments the retained lines use; the article's own Elsevier class file and its \texttt{amssym.def} input are neither shipped nor input, and the retained lines use no command that only those files provide; sequential numbering of the retained environments in order of appearance, whereas the article prints the counted lemma as Lemma 3; equation numbering is not reproduced and every cross-reference inside the retained text resolves to the wrapper's numbering. No citation occurs inside any retained line, so the wrapper carries no bibliography. No asset-inclusion command, no drawing environment and no image asset occurs in this package.


\section{The counted unit: Lemma 3 and its complete original proof}

\begin{lemma}
\label{lemma: limit} As well as C1) - C3) assume that the parameter space $%
\Theta $ is a compact subset of $\mathbb{R}^{p}$. Then with $\Psi _{n}\left( 
\boldsymbol{\theta }\right) \overset{def}{=}n^{-1}\sum_{i,j}\psi _{\sigma
}\left( Y_{j}\left( \boldsymbol{x}_{i}\right) -\boldsymbol{f}^{\T
}\left( \boldsymbol{x}_{i}\right) \left( \boldsymbol{\theta }+\boldsymbol{%
\tilde{\theta}}_{n}\right) \right) \boldsymbol{f}\left( \boldsymbol{x}%
_{i}\right) ,$ the function 
\begin{equation*}
\Psi \left( \boldsymbol{\theta }\right) =\int_{\mathcal{\chi }}E\left[ \psi
_{\sigma }\left( Y\left( \boldsymbol{x}\right) -\boldsymbol{f}^{\T
}\left( \boldsymbol{x}\right) \left( \boldsymbol{\theta }+\boldsymbol{\theta 
}_{0}\right) \right) \right] \boldsymbol{f}\left( \boldsymbol{x}\right) \xi
_{\ast }\left( d\boldsymbol{x}\right)
\end{equation*}%
satisfies 
\begin{eqnarray*}
(i) &&\sup_{\boldsymbol{\theta }\in \boldsymbol{\Theta }}\left\Vert \Psi _{n}\left( \boldsymbol{%
\theta }\right) -\Psi \left( \boldsymbol{\theta }\right) \right\Vert \overset%
{pr}{\rightarrow }0, \\
(ii) &\text{ }&\text{for every }\delta >0\text{, }\inf_{\left\Vert 
\boldsymbol{\theta }\right\Vert \geq \delta }\left\Vert \Psi \left( 
\boldsymbol{\theta }\right) \right\Vert >0=\left\Vert\Psi ( \boldsymbol{0}_p\right)\Vert
.
\end{eqnarray*}
\end{lemma}

\subsection{Proof of Lemma \protect\ref{lemma: limit}}

With $\boldsymbol{\tilde{\Psi}}_{n}(\boldsymbol{\theta })\overset{def}{=}%
\int_{\chi }E\!\left[ \psi _{\sigma }\!\left( Y(\boldsymbol{x})-\boldsymbol{f%
}^{\T }(\boldsymbol{x})(\boldsymbol{\theta }+\widetilde{\boldsymbol{%
\theta }}_{n})\right) \right] \boldsymbol{f}(\boldsymbol{x})\,\xi _{n}(d%
\boldsymbol{x})=E\left[ \Psi _{n}\left( \boldsymbol{\theta }\right) \right] $
we shall establish (i) of the Lemma by showing:

\begin{enumerate}
\item[(a)] $\sup_{\boldsymbol{\theta }\in \boldsymbol{\Theta }}\left\Vert 
\boldsymbol{\tilde{\Psi}}_{n}\left( \boldsymbol{\theta }\right) -\boldsymbol{%
\Psi }\left( \boldsymbol{\theta }\right) \right\Vert \rightarrow 0$,

\item[(b)] $\sup_{\boldsymbol{\theta }\in \boldsymbol{\Theta }}\left\Vert 
\boldsymbol{\Psi }_{n}\left( \boldsymbol{\theta }\right) -\boldsymbol{\tilde{%
\Psi}}_{n}\left( \boldsymbol{\theta }\right) \right\Vert \overset{pr}{%
\rightarrow }0$.
\end{enumerate}

For (a), with $\boldsymbol{h}(\boldsymbol{x},\boldsymbol{\theta },%
\boldsymbol{\eta })\overset{def}{=}E\!\left[ \psi _{\sigma }\!\left( Y(%
\boldsymbol{x})-\boldsymbol{f}^{\T }(\boldsymbol{x})(\boldsymbol{\theta }%
+\boldsymbol{\eta })\right) \right] \boldsymbol{f}(\boldsymbol{x})$ we have $%
\boldsymbol{\Psi }(\boldsymbol{\theta })=\int_{\chi }\boldsymbol{h}(%
\boldsymbol{x},\boldsymbol{\theta },\boldsymbol{\theta }_{0})\,\xi _{\ast }(d%
\boldsymbol{x})$ and $\boldsymbol{\tilde{\Psi}}_{n}\left( \boldsymbol{\theta 
}\right) =\int_{\chi }\boldsymbol{h}(\boldsymbol{x},\boldsymbol{\theta },%
\boldsymbol{\tilde{\theta}}_{n})\,\xi _{n}(d\boldsymbol{x})$, with 
\begin{eqnarray*}
\boldsymbol{\Psi }(\boldsymbol{\theta })-\boldsymbol{\tilde{\Psi}}_{n}\left( 
\boldsymbol{\theta }\right) &=&\int_{\chi }\left( \boldsymbol{h}(\boldsymbol{%
x},\boldsymbol{\theta },\boldsymbol{\theta }_{0})-\boldsymbol{h}(\boldsymbol{%
x},\boldsymbol{\theta },\boldsymbol{\tilde{\theta}}_{n})\right) \,\xi _{n}(d%
\boldsymbol{x})+\int_{\chi }\boldsymbol{h}(\boldsymbol{x},\boldsymbol{\theta 
},\boldsymbol{\theta }_{0})\,(\xi _{\ast }-\xi _{n})(d\boldsymbol{x}) \\
&=&\boldsymbol{\alpha }_{n}(\boldsymbol{\theta })+\boldsymbol{\beta }_{n}(%
\boldsymbol{\theta })\text{, say.}
\end{eqnarray*}%
Since $\boldsymbol{h}(\boldsymbol{x},\boldsymbol{\theta },\boldsymbol{\eta }%
) $ is differentiable w.r.t. $\boldsymbol{\eta }$, and both $\psi _{\sigma
}^{\prime }$ and $\boldsymbol{f}$ are bounded, there is $C>0$ for which $%
\left\Vert \boldsymbol{h}(\boldsymbol{x},\boldsymbol{\theta },\boldsymbol{%
\theta }_{0})-\boldsymbol{h}(\boldsymbol{x},\boldsymbol{\theta },\boldsymbol{%
\tilde{\theta}}_{n})\right\Vert \leq C\left\Vert \boldsymbol{\theta }_{0}-%
\boldsymbol{\tilde{\theta}}_{n}\right\Vert $, uniformly in $\boldsymbol{x}%
\in \chi $ and $\boldsymbol{\theta }\in \Theta $. Hence 
\begin{equation}
\sup_{\boldsymbol{\theta }\in \Theta }\left\Vert \boldsymbol{\alpha }_{n}(%
\boldsymbol{\theta })\right\Vert \leq C\left\Vert \boldsymbol{\theta }_{0}-%
\boldsymbol{\tilde{\theta}}_{n}\right\Vert \rightarrow 0.  \label{first sup}
\end{equation}%
With $\boldsymbol{g}_{\boldsymbol{\theta }}(\boldsymbol{x})\overset{def}{=}%
\boldsymbol{h}(\boldsymbol{x},\boldsymbol{\theta },\boldsymbol{\theta }_{0})$
we obtain, in a similar fashion, $\sup_{\boldsymbol{x}\in \chi }\Vert 
\boldsymbol{g}_{\boldsymbol{\theta }_{1}}(\boldsymbol{x})-\boldsymbol{g}_{%
\boldsymbol{\theta }_{2}}(\boldsymbol{x})\Vert \leq C\Vert \boldsymbol{%
\theta }_{1}-\boldsymbol{\theta }_{2}\Vert $. Then since $\boldsymbol{\Theta 
}$ is compact, for every $\delta >0$ there exist $\boldsymbol{\theta }%
^{(1)},\dots ,\boldsymbol{\theta }^{(K)}\in \Theta $ such that every $%
\boldsymbol{\theta }\in \boldsymbol{\Theta }$ lies within distance $\delta $
of one of these points. Therefore, 
\begin{equation*}
\sup_{\boldsymbol{\theta }\in \Theta }\left\Vert \boldsymbol{\beta }_{n}(%
\boldsymbol{\theta })\right\Vert \leq 2C\delta +\max_{1\leq k\leq
K}\left\Vert \int_{\chi }\boldsymbol{g}_{\boldsymbol{\theta }^{(k)}}(%
\boldsymbol{x})\,\xi _{\ast }(d\boldsymbol{x})-\int_{\chi }\boldsymbol{g}_{%
\boldsymbol{\theta }^{(k)}}(\boldsymbol{x})\,\xi _{n}(d\boldsymbol{x}%
)\right\Vert .
\end{equation*}%
That $\xi _{n}\overset{d}{\rightarrow }\xi _{\ast }$ implies that the
preceding norms, hence their maximum, converge to $0$. Letting $\delta
\rightarrow 0$, we conclude that $\sup_{\boldsymbol{\theta }\in \Theta
}\left\Vert \boldsymbol{\beta }_{n}(\boldsymbol{\theta })\right\Vert
\rightarrow 0$, which together with (\ref{first sup}) yields (a).

For (b), we write $\boldsymbol{\Delta }_{n}(\boldsymbol{\theta })\overset{def%
}{=}\boldsymbol{\Psi }_{n}(\boldsymbol{\theta })-\boldsymbol{\tilde{\Psi}}%
_{n}\left( \boldsymbol{\theta }\right) \ =n^{-1}\sum_{i,j}\left\{ 
\boldsymbol{\zeta }_{ij,n}(\boldsymbol{\theta })-E[\boldsymbol{\zeta }%
_{ij,n}(\boldsymbol{\theta })]\right\} ,$ where 
\begin{equation*}
\boldsymbol{\zeta }_{ij,n}(\boldsymbol{\theta })=\psi _{\sigma }\!\left(
Y_{j}(\boldsymbol{x}_{i})-\boldsymbol{f}^{\T }\left( \boldsymbol{x}%
_{i}\right) (\boldsymbol{\theta }+\widetilde{\boldsymbol{\theta }}%
_{n})\right) \boldsymbol{f}(\boldsymbol{x}_{i}).
\end{equation*}%
Arguing as above, there exists $C>0$ such that for all $\boldsymbol{\theta }%
_{1},\boldsymbol{\theta }_{2}\in \boldsymbol{\Theta }$ and each pair $(i,j)$%
, 
\begin{equation*}
\left\Vert \boldsymbol{\zeta }_{ij,n}(\boldsymbol{\theta }_{1})-\boldsymbol{%
\zeta }_{ij,n}(\boldsymbol{\theta }_{2})\right\Vert \leq C\left\Vert 
\boldsymbol{\theta }_{1}-\boldsymbol{\theta }_{2}\right\Vert .
\end{equation*}%
Taking expectations gives the same bound for the centred version, so that 
\begin{equation*}
\left\Vert \boldsymbol{\Delta }_{n}(\boldsymbol{\theta }_{1})-\boldsymbol{%
\Delta }_{n}(\boldsymbol{\theta }_{2})\right\Vert \leq 2C\left\Vert 
\boldsymbol{\theta }_{1}-\boldsymbol{\theta }_{2}\right\Vert .
\end{equation*}%
Thus $\Delta _{n}$ is uniformly Lipschitz on $\boldsymbol{\Theta }$, with a
constant independent of $n$.

Fix $\varepsilon >0$. Choose $\delta >0$ such that $2C\delta <\varepsilon /2$%
, and let $\boldsymbol{\theta }^{(1)},\dots ,\boldsymbol{\theta }^{(K)}$ be
a finite $\delta $-net of $\Theta $, as above. Then 
\begin{equation*}
\sup_{\boldsymbol{\theta }\in \Theta }\left\Vert \boldsymbol{\Delta }_{n}(%
\boldsymbol{\theta })\right\Vert \leq \max_{1\leq k\leq K}\left\Vert 
\boldsymbol{\Delta }_{n}(\boldsymbol{\theta }^{(k)})\right\Vert +\varepsilon
/2,
\end{equation*}%
and it remains only to show that 
\begin{equation}
\max_{1\leq k\leq K}\left\Vert \boldsymbol{\Delta }_{n}(\boldsymbol{\theta }%
^{(k)})\right\Vert \overset{pr}{\rightarrow }0.  \label{remains}
\end{equation}%
Write $\Delta _{n}(\boldsymbol{\theta }^{(k)})=n^{-1}\sum_{i,j}\left( 
\boldsymbol{\zeta }_{ij,n}(\boldsymbol{\theta }^{(k)})-E\left[ \boldsymbol{%
\zeta }_{ij,n}(\boldsymbol{\theta }^{(k)})\right] \right) =n^{-1}\sum_{i,j}%
\boldsymbol{Z}_{ij,n}^{(k)}$, say. We can write $\boldsymbol{\zeta }_{ij,n}(%
\boldsymbol{\theta })$ as $\boldsymbol{\zeta }_{ij,n}(\boldsymbol{\theta })=\psi _{\sigma }\!\left(
\varepsilon _{ij}+e_{i,n}\left( \boldsymbol{\theta }\right) \right) 
\boldsymbol{f}(\boldsymbol{x}_{i})
$,
for certain functions $e_{i,n}\left( \boldsymbol{\theta }\right) $,
continuous on $\boldsymbol{\Theta }$ hence uniformly bounded there. Then $E[%
\boldsymbol{Z}_{ij,n}^{(k)}]=\boldsymbol{0}$ and{\ there is }$C>0$ for which 
{$\sup_{i,j,n,k}E\Vert \boldsymbol{Z}_{ij,n}^{(k)}\Vert ^{2}\leq C$}. Since
the $\boldsymbol{Z}_{ij,n}^{(k)}$ are independent for each $(k,n)$, \textsc{%
var}$\left[ \Delta _{n}(\boldsymbol{\theta }^{(k)})\right] \leq
n^{-1}C\rightarrow 0$, implying by Chebyshev's inequality that $\boldsymbol{%
\Delta }_{n}(\boldsymbol{\theta }^{(k)})\overset{pr}{\rightarrow }0$ for
each $k=1,\dots ,K$. Now (\ref{remains}) follows, thus completing the proof
of (b).

For part (ii) of the Lemma, recall that the design space
is finite and that each $\xi _{n}\left( \boldsymbol{x}_{i}\right) $ tends to 
$\xi _{\ast }\left( \boldsymbol{x}_{i}\right) $, so that
\begin{equation}
\int_{\mathcal{\chi }}E\left[ \psi _{\sigma }\left( Y\left( \boldsymbol{x}%
\right) -\boldsymbol{f}^{\T }\left( \boldsymbol{x}\right) \boldsymbol{%
\theta }_{0}\right) \right] \boldsymbol{f}\left( \boldsymbol{x}\right) \xi
_{\ast }\left( d\boldsymbol{x}\right) =\lim_{n\rightarrow \infty }\int_{%
\mathcal{\chi }}E\left[ \psi _{\sigma }\left( \varepsilon +\tau \left( 
\boldsymbol{x}\right) \right) \right] \boldsymbol{f}\left( \boldsymbol{x}%
\right) \xi _{n}\left( d\boldsymbol{x}\right) =\boldsymbol{0}_p,  \label{root}
\end{equation}%
since $\tau \left( \boldsymbol{x}\right) $ is $O\left( n^{-1/2}\right) $. It
follows that $\Psi \left( \boldsymbol{0}_p\right) =\boldsymbol{0}_p $.

By (\ref{root}), $\boldsymbol{\theta }_{0}$ is a stationary point of the
function $k\left( \boldsymbol{\theta }\right) \overset{def}{=}\int_{\mathcal{%
\chi }}E\left[ \rho _{\sigma }\left( Y\left( \boldsymbol{x}\right) -%
\boldsymbol{f}^{\T }\left( \boldsymbol{x}\right) \boldsymbol{\theta }%
\right) \right] \xi _{\ast }\left( d\boldsymbol{x}\right) $, where $\rho
_{\sigma }(x)=\int^{x}\psi _{\sigma }\left( t\right) dt$. The Hessian is $%
\ddot{k}\left( \boldsymbol{\theta }\right) =\int_{\mathcal{\chi }}E\left[
\psi _{\sigma }^{\prime }\left( \varepsilon +M_{\boldsymbol{\theta }}\left( 
\boldsymbol{x}\right) \right) \right] \boldsymbol{f}\left( \boldsymbol{x}%
\right) \boldsymbol{f}^{\T }\left( \boldsymbol{x}\right) \xi _{\ast
}\left( d\boldsymbol{x}\right) ,$ for $M_{\boldsymbol{\theta }}\left( 
\boldsymbol{x}\right) =E\left[  Y\left( \boldsymbol{x}%
\right)] -\boldsymbol{f} ^{\T }\left( \boldsymbol{x}\right) \boldsymbol{\theta }\right) $. Since $\psi_{\sigma }$ is weakly increasing, $\psi _{\sigma }^{\prime }\left(
\varepsilon +M_{\boldsymbol{\theta }}\left( \boldsymbol{x}\right) \right)
\geq 0$ and so $E\left[ \psi _{\sigma }^{\prime }\left( \varepsilon +M_{%
\boldsymbol{\theta }}\left( \boldsymbol{x}\right) \right) \right] $ must be
strictly positive, else $\psi _{\sigma }^{\prime }\left( \varepsilon +M_{%
\boldsymbol{\theta }}\left( \boldsymbol{x}\right) \right) \equiv 0$ (a.s.).
Put 
\begin{equation*}
c\left( \boldsymbol{\theta }\right) =\min_{i=1,...,N}\left\{ E\left[ \psi
_{\sigma }^{\prime }\left( \varepsilon +M_{\boldsymbol{\theta }}\left( 
\boldsymbol{x}_{i}\right) \right) \right] \left\vert {}\right. \xi _{\ast
}\left( \boldsymbol{x}_{i}\right) >0\right\} .
\end{equation*}%
Then $c\left( \boldsymbol{\theta }\right) >0$ and $\ddot{k}\left( 
\boldsymbol{\theta }\right) \succeq c\left( \boldsymbol{\theta }\right)
\int_{\mathcal{\chi }}\boldsymbol{f}\left( \boldsymbol{x}\right) \boldsymbol{%
f}^{\T }\left( \boldsymbol{x}\right) \xi _{\ast }\left( d\boldsymbol{x}%
\right) \succ \boldsymbol{0},$ so that $\boldsymbol{\theta }_{0}$ is the
unique stationary point of the strictly convex function $k\left( \boldsymbol{%
\theta }\right) $ and $\boldsymbol{0}_p$ is the unique zero of $\left\Vert 
\boldsymbol{\Psi }\left( \boldsymbol{\theta }\right) \right\Vert $.

Since $\boldsymbol{\Theta }$ is compact, its closed subset $\boldsymbol{%
\Theta }_{\delta }$ defined by $\left\Vert \boldsymbol{\theta }\right\Vert
\geq \delta $ is also compact, hence the continuous function $\left\Vert 
\boldsymbol{\Psi }\left( \boldsymbol{\theta }\right) \right\Vert $ attains
its $\inf $ there:%
\begin{equation*}
\inf_{\boldsymbol{\Theta }_{\delta }}\left\Vert \boldsymbol{\Psi }\left( 
\boldsymbol{\theta }\right) \right\Vert =\min_{\boldsymbol{\Theta }_{\delta
}}\left\Vert \boldsymbol{\Psi }\left( \boldsymbol{\theta }\right)
\right\Vert >0.
\end{equation*}%
This proves (ii), and completes the proof of Lemma \ref{lemma: limit}%
.\hfill\ $\square $

\medskip


\end{document}
